\section{Materiali e metodi}

\subsection{Algoritmo di Shor, validazione e definizione del successo}\label{sec:implementazione}

Scelta una base \(a\) coprima con il numero \(N\) da fattorizzare, la successione
\(a^{x} \bmod N\) è periodica; il periodo, detto ordine \(r\), è il minimo intero
positivo per cui \(a^{r} \equiv 1 \pmod{N}\). La parte quantistica lo stima con stima
di fase e QFT inversa; la parte classica ricava i fattori per approssimazione razionale
e massimo comune divisore, e li verifica.

L'istanza principale è \(N = 15\), \(a = 7\), con otto qubit di controllo e quattro di
lavoro; l'ordine è \(r = 4\) e gli esiti ideali del registro di controllo sono 0, 64,
128 e 192. Ogni esecuzione con misura, detta \emph{shot}, restituisce un solo esito;
molti shot formano l'istogramma su cui opera il post-processing. Il successo è il
recupero di fattori non banali, non la ricostruzione esatta dell'ordine. Vi conducono
tre dei quattro picchi ideali, con probabilità complessiva del 75\%. La verifica
classica è però permissiva e accetta 63 dei 256 valori, tanto che esiti del tutto
uniformi darebbero comunque un successo di \(63/256 \simeq 24{,}61\%\). Questo riferimento dipende dall'istanza e dal post-processing e non misura la fedeltà dello stato quantistico. Mostra inoltre che fattorizzazioni corrette possono essere ottenute anche quando la distribuzione degli esiti di misura del registro di controllo ha perso gran parte della struttura ideale dei picchi \cite{smolin2013}.

La correttezza dei circuiti è verificata indipendentemente dal successo della fattorizzazione,
controllando operatori modulari, ripristino dei qubit ausiliari e accordo con le
distribuzioni teoriche. Oltre a \(N = 15\), si considerano le istanze \(N = 21\), \(a = 2\) e
\(N = 35\), \(a = 6\), che, validate idealmente, usano
l'aritmetica reversibile di Beauregard \cite{beauregard2002}; solo \(N = 21\) entra nelle
prove rumorose esplorative sull'AQFT, perché per \(N = 35\) le risorse hardware disponibili
non consentivano la simulazione di un circuito di Shor adeguato.
Si impiegano Qiskit/Aer, scikit-learn, Stim, PyMatching e ldpc, registrando versioni,
parametri, semi casuali e identificativi dei circuiti.

\subsection{Sensibilità di Shor agli errori di porta}

Il punto di partenza dello studio è stato misurare quanto l'algoritmo di Shor sia robusto
agli errori di porta. Nel circuito dell'istanza principale, compilato nelle porte native
RZ, SX, X e CX (profondità 412, con 224 porte CX e 70 SX), si introducono errori di Pauli
indipendenti dopo ogni porta, escluse le rotazioni RZ, che sono virtuali. L'intensità del
rumore è fissata dalla probabilità \(p_g\) che dopo una porta si verifichi un errore di
Pauli diverso dall'identità, fatta variare su 13 valori fra 0 e 0,5.

Il modello di rumore è fenomenologico e uniforme e non riproduce la calibrazione di un
dispositivo reale. Per questo \(p_g\) va inteso come probabilità di errore per singola
porta: non è direttamente confrontabile con la probabilità di fallimento logico \(p_L\)
dei codici di correzione, perché l'algoritmo di Shor non viene compilato in forma
fault-tolerant.

\subsection{Mitigazione degli errori mediante post-processing e machine learning}

La prima strategia agisce a valle della misura: senza modificare il circuito, cerca di
ricavare i fattori dagli esiti rumorosi scegliendo meglio quali verificare, e valuta se un
modello di machine learning renda questa scelta più efficace. Tre strategie di selezione,
con estrattore dei fattori e regola dei pareggi comuni, operano sugli stessi istogrammi,
uno per iterazione quantistica, da 1.024 shot ciascuno. TOP-1 (Metodo~1) verifica il solo
esito più frequente; TOP-4 fino a quattro candidati ordinati per frequenza; il Metodo~2
applica TOP-4 solo se un classificatore predice il successo del candidato dominante. Il
confronto fra Metodo~2 e TOP-4 senza filtro è lo studio di ablazione, che stabilisce se il
vantaggio venga dal modello appreso o dal semplice provare più candidati.

Due scenari illustrativi, UC1 di riferimento e UC2 di stress, includono
depolarizzazione, rilassamento termico ed errori di lettura; per ciascuno scenario, 2.000
istogrammi con parametri variati entro \(\pm50\%\) del nominale sono divisi
60\%/20\%/20\% fra addestramento, selezione e verifica. Random forest, support vector
machine (SVM) e perceptron multistrato sono confrontati su F1; il modello scelto è
valutato su 400 istogrammi indipendenti per scenario. La metrica primaria è il numero di
iterazioni per ottenere i fattori; l'accuratezza è stata considerata un elemento
secondario nello studio del classificatore, mentre il tempo di calcolo è stato rapportato
alle risorse hardware disponibili. Si usano 30 repliche appaiate per scenario e al massimo
50 iterazioni.

\subsection{Correzione quantistica degli errori: surface code e codici qLDPC}

La seconda strategia protegge l'informazione prima che il rumore la distrugga. Nella
correzione quantistica degli errori, l'informazione di un qubit logico viene distribuita su
più qubit fisici; le sindromi permettono di rilevare gli errori senza misurare
direttamente lo stato logico e guidano il decoder nella scelta della correzione. Le
prestazioni sono misurate attraverso la probabilità di fallimento logico \(p_L\).

Codice a ripetizione e codice di Steane sono utilizzati come verifiche preliminari, mentre
l'analisi principale riguarda circuiti di memoria basati su surface code
\cite{steane1996,fowler2012,googleQECBelowThreshold2025}, con errori nelle operazioni,
nella preparazione e nella misura e, in alcune configurazioni, crosstalk tra qubit
adiacenti. Il decoder di riferimento è il \emph{minimum-weight perfect matching} (MWPM),
che associa a coppie le sindromi attivate scegliendo la configurazione di errori più
probabile. MWPM è confrontato con una rete neurale autonoma e con una strategia ibrida, la
cui soglia d'intervento è selezionata sul validation set e mantenuta fissa sul test set
\cite{alphaqubit2024}.

Uno studio esplorativo confronta inoltre il surface code con il Gross code di IBM, un
codice quantum Low-Density Parity-Check (qLDPC) che codifica 12 qubit logici in 144 qubit
di dato \cite{bravyi2024}, nel modello code-capacity, con errori sui soli qubit di dato e
sindromi perfette. BP+OSD (\emph{belief propagation with ordered statistics decoding}) è
usato come ulteriore decoder di confronto sul surface code e come decoder del Gross code
\cite{roffe2020}.

Il confronto mostra che i due approcci non sono equivalenti: a parità di qubit logici il
Gross code richiede molti meno qubit di dato, con un vantaggio che dipende in modo decisivo
dal decoder (Sezione~\ref{sec:risultati-qec}).

\subsection{Riduzione del rumore con la AQFT}

La terza strategia riduce l'effetto del rumore alla fonte, diminuendo il numero di porte su
cui può agire. La QFT approssimata (AQFT) elimina le rotazioni controllate di piccolo angolo
della QFT, che contribuiscono poco all'informazione di fase: il circuito perde un po' di
precisione, ma ha meno porte e minore profondità, e quindi meno occasioni di errore
\cite{barencoAQFT1996}. La domanda è se il rumore evitato compensi la precisione persa.

Si analizzano separatamente il troncamento della QFT inversa finale di Shor per \(N = 15\) e
quello delle trasformate impiegate nell'aritmetica modulare per \(N = 21\); benchmark
distinti di stima di fase, con registri da sei a dieci qubit, isolano l'effetto del
troncamento. Il modello di rumore è costruito a partire dalle infedeltà delle porte e dagli
errori di lettura della calibrazione offline FakeSherbrooke, senza aggiungere contributi
\(T_1/T_2\), mantenendo invariati layout e semi casuali. Il grado di troncamento viene
selezionato su un insieme dedicato e valutato su dati di verifica indipendenti; gli shot per
configurazione, 8.192 nel pilota per \(N = 15\) e 256 per \(N = 21\), sono suddivisi
equamente tra selezione e verifica.
